% ======================================================================
% SECTION 5: HR 9505 - Real-Time Patient-Sponsor Direct Communication
% Act of 2026
% Revision of FDA Real-Time Clinical Trials Press Announcement April
% 2026 and FDA Decentralized Clinical Trials Guidance 2024.
% ======================================================================

\hspace{0.5cm} HR 9505 is a \textit{Revision of the FDA Real-Time Clinical Trials
Press Announcement of April 28, 2026 and the FDA Decentralized
Clinical Trials Final Guidance of 2024}, drawing the chain of prior
authority from the FDA's own validation of the Paradigm-Health-style
real-time signal-sharing conduit through the FDA decentralized-
elements final guidance \cite{fda_rtct_press_2026,
fda_dct_guidance_2024, fda_dhts_drug_development_2026}. The bill
extends a sponsor-initiated proof of concept (TRAVERSE Phase 2 in
mantle cell lymphoma at MD Anderson and Penn; STREAM-SCLC Phase 1b
at Mass General) into a patient-initiated direct sponsor
communication right, while remaining respectful of the FDA's
leadership posture as the agency the bill builds on rather than
against. The bill implements the fifth dimension of the operational
definition of patient control: toxicity warnings and adverse-event
self-management through ePRO, ASyMS, and eRAPID infrastructure
\cite{rocque2025remote_symptom_monitoring, kearney2009asyms,
maguire2017esmart, absolom2021erapid}.

\subsection{Findings and Declarations}

\hspace{0.5cm} (a) FDA Commissioner Marty Makary, in the April 28, 2026 press
announcement, stated that "for 60 years, we've been conducting
clinical trials in the same way" and framed real-time clinical
trials as a transformative shift in how trials are run. FDA Chief
AI Officer Jeremy Walsh emphasized the transformative potential of
real-time trials in the same announcement
\cite{fda_rtct_press_2026}.

(b) Two announced proof-of-concept trials anchor the FDA's
validation: AstraZeneca's TRAVERSE Phase 2 trial in mantle cell
lymphoma at MD Anderson and the University of Pennsylvania, and
Amgen's STREAM-SCLC Phase 1b trial at Massachusetts General
Hospital \cite{fda_rtct_press_2026}.

(c) Paradigm Health is the validated technical conduit for real-time signal sharing in both proof-of-concept trials
\cite{fda_rtct_press_2026}.

(d) The author's prior 168-hour autonomous sponsor simulation at repository file location
\texttt{sponsor/final\_paper/168\_hours/} demonstrated that hourly
commits from a sponsor-side server (24 repository commits per day, 168 commits
across seven days) sustain a reduced size trial-data flow on minimal
commodity hardware (Core i5-6200U with 4 GB RAM, see
\texttt{sponsor/final\_paper/168\_hours/instructions/core\_i5\_6200u\_4gb/README.md})
\cite{kawchak_2026_19396256}.

(e) The ePRO and remote symptom monitoring evidence (Rocque 2025;
ASyMS; eRAPID) supports the bill without overclaiming
\cite{rocque2025remote_symptom_monitoring, kearney2009asyms,
maguire2017esmart, absolom2021erapid}.

(f) The FDA pilot is sponsor-initiated; this bill extends the same
real-time channel to patient-initiated communication while
preserving the FDA's regulatory leadership.

\subsection{Definitions}

\hspace{0.5cm} (a) \billterm{Patient-sponsor direct channel} means a HIPAA-compliant, FHIR-based, Paradigm-Health-style aggregator endpoint
that the patient may use to send the patient's own structured data,
modification requests, and adverse-event reports directly to the
trial sponsor's safety pharmacovigilance system.\ The operational
baseline is the FastAPI sponsor control server at
\texttt{sponsor/final\_paper/scripts/sponsor\_server/} and the
safety-pharmacovigilance section of the prior sponsor paper at
\texttt{sponsor/paper/sections/safety\_pharmacovigilance.tex}
\cite{kawchak_2026_19396256}.

(b) \billterm{Real-time} means data from the patient's qualified
Physical AI executor, ePRO endpoint, or wearable reaches the
sponsor within one hour of generation. The hourly commit cadence
documented in the four-simulation paper at
\texttt{new-trial/national-24-7-trial/paper/sections/results.tex}
(Simulation 1: 24 commits per day; 84 commits across 84 hours) and
the 168-hour extension at \texttt{sponsor/final\_paper/168\_hours/}
(Simulation 4: 168 GitHub commits across seven days) provide the
operational baseline \cite{kawchak_2026_19994945,
kawchak_2026_19396256}.

(c) \billterm{Sponsor-side patient interface} means the cancer patient
facing endpoint that mirrors
\texttt{sponsor/final\_paper/scripts/sponsor\_server/} but is
accessible to the patient under HIPAA-compliant authentication.

\subsection{Patient-Sponsor Direct Communication Rights}

\hspace{0.5cm} (a) Right to send structured adverse-event reports (Common Toxicity
Criteria graded) directly to the sponsor's pharmacovigilance system
within one hour, drawing on the ePRO infrastructure documented in
Rocque 2025 and the ASyMS / eRAPID systems
\cite{rocque2025remote_symptom_monitoring, kearney2009asyms,
maguire2017esmart, absolom2021erapid}.

(b) Right to receive sponsor-side acknowledgment of every patient
submission within one hour, with the acknowledgment carrying a
hash-chained provenance record per the MCP-PAI standard documented
at \texttt{national-platform/national\_mcp/}
\cite{kawchak_2026_19244918}.

(c) Right to receive a per-cycle digital twin update reflecting
the patient's own data submissions.\ The technical baseline is the
digital-twin initialization Python module at repository file location
\texttt{patient-journey/stage\_03\_digital\_twin.py} and the
patient-modeling framework at
\texttt{digital-twins/patient-modeling/}
\cite{kawchak_2026_19119939}.

(d) Right to escalate any sponsor-side delay beyond one hour to
the FDA RTCT pilot oversight team, mirroring the escalation
pattern at \texttt{sponsor/final\_paper/scripts/coordination/}
\cite{kawchak_2026_19396256}.

(e) Right to receive a real-time read-back of how the patient's
data is being used in the sponsor's safety analyses, with the
privacy envelope governed by 45 CFR 46 and the data-protection
framework in HR 9507 \cite{hhs_45cfr46_protection_human_subjects}.

(f) Right to opt out of any individual sponsor-data exchange while
remaining enrolled in the trial. The patient retains discretion
over which fields the sponsor sees and may withdraw any individual
field without withdrawing from the trial.

\subsection{Implementation}

\hspace{0.5cm} (a) The Secretary of Health and Human Services, acting through the
Food and Drug Administration, the Office of the National
Coordinator for Health Information Technology, and the Department
of Health and Human Services, shall require all FDA RTCT pilot
trials to expose a patient-sponsor direct channel within 24 months
of enactment.

(b) Implementation aligns with the FDA's own RTCT pilot program and
the FDA Oncology Center of Excellence Advancing Oncology
Decentralized Trials directory
\cite{fda_oce_advancing_oncology_dct_2024}.

(c) The NIH funding opportunity for digital-health-technology-derived endpoints in cancer trials (PAR-25-170) supplies the
research-funding bridge to expand the patient-sponsor direct
channel beyond the initial pilot
\cite{nih_par25_170_dht_endpoints}.

\subsection{Reporting and Enforcement}

\hspace{0.5cm} (a) Annual federal report on patient-sponsor channel uptake and
median sponsor acknowledgment latency. The federal Patient Control
Dashboard's Domain 3 (Control-in-Action) subset anchors this
report: alert response time, false-alert burden, median sponsor
acknowledgment latency, and serious adverse events detected via
patient-initiated channel versus sponsor-initiated channel
\cite{rocque2025remote_symptom_monitoring}.

(b) Civil penalties for sponsors that fail to acknowledge within
one hour.

(c) 
\texttt{sponsor/final\_paper/scripts/core\_agents/audit\_trail\_manager.py} audit-trail provenance model
provides the technical baseline for hash-chained provenance records
\cite{kawchak_2026_19396256}.

\subsection{Discussion and Limitation}

\hspace{0.5cm} The ePRO and remote symptom monitoring evidence supports the proposed new bill
without overclaiming.\ The author's prior simulation evidence (consisting of the
four LLM simulations in
\texttt{new-trial/national-24-7-trial/paper/full-paper/}, the
single-patient journey at \texttt{patient-journey/}, and the 24-
hour and 168-hour sponsor simulations in \texttt{sponsor/final\_paper/})
supplies the operational substrate that the bill regulates.\ The
FDA's "demonstrated proof-of-concept" framing (TRAVERSE, STREAM-
SCLC, Paradigm Health) is the regulatory opening that HR 9505
builds on \cite{kawchak_2026_19994945, kawchak_2026_19119939,
kawchak_2026_19396256, fda_rtct_press_2026,
fda_dhts_drug_development_2026}. The principal limitation is that
Paradigm-Health-style aggregator infrastructure has not yet been
demonstrated to generalize beyond TRAVERSE and STREAM-SCLC; the
FDA RTCT pilot has not released aggregate metrics; and the 24-month
deadline is aggressive \cite{fda_rtct_press_2026}.

\subsection{Patient-Control Framings}

\hspace{0.5cm} For the individual patient, a 58-year-old female NSCLC patient
matching PAT-2026-0042 sends a Grade 2 fatigue adverse-event report
from her home wearable on day 14 of cycle 22 of her 35
pembrolizumab cycles.\ The sponsor's safety pharmacovigilance system
acknowledges the report within 12 minutes and updates her per-cycle
digital twin within the same hour. The patient reviews the read-
back showing how her adverse-event data informed the cumulative
safety analysis at
\texttt{sponsor/final\_paper/scripts/output/sponsor\_24h\_summary.json}
\cite{kawchak_2026_19119939, kawchak_2026_19396256}.

For broad adoption, all FDA RTCT pilot trials must expose a
patient-sponsor direct channel within 24 months. The bill positions
the United States as Number 1 in the world for patient-initiated
real-time trial communication, building respectfully on the FDA's
own April 28, 2026 RTCT framework and the 168-hour autonomous
sponsor simulation in \texttt{sponsor/final\_paper/168\_hours/}
\cite{kawchak_2026_19396256, kawchak_2026_19994945,
fda_rtct_press_2026}.
